\chapter{Estação de Monitoramento IoT}
\label{cap:projeto-final}

Este capítulo final reúne, em um único projeto, praticamente tudo o que foi apresentado nesta apostila: tipos de dados, estruturas de controle, funções, vetores, \code{structs} e \code{enums}, além de GPIO, ADC, comunicação serial e interrupções. O objetivo é consolidar esses conceitos através de um firmware completo e funcional — não apenas um exemplo isolado de cada periférico.

\section{Especificação do projeto}

Uma \textbf{estação de monitoramento} de luminosidade ambiente: o \ESP{} lê continuamente um sensor LDR, mantém uma média móvel das últimas leituras para suavizar ruído, registra o estado do ambiente pela porta serial e acende um LED de alerta sempre que o ambiente estiver escuro. Um botão alterna entre modo \textbf{normal} (só registra alertas) e modo \textbf{verboso} (registra cada leitura) — sem nunca bloquear o laço principal com \code{delay}.

\subsection{Materiais}

\begin{itemize}
    \item 1 placa de desenvolvimento \ESP{};
    \item 1 LED (com resistor limitador em série, \textasciitilde220~$\Omega$), como indicador de alerta;
    \item 1 sensor LDR, em um divisor de tensão com um resistor fixo (\textasciitilde10~k$\Omega$);
    \item 1 botão \emph{push-button};
    \item Protoboard e jumpers.
\end{itemize}

\subsection{Ligações sugeridas}

\begin{itemize}
    \item LED de alerta $\rightarrow$ GPIO2 (através do resistor limitador, ao GND);
    \item Ponto médio do divisor de tensão do LDR $\rightarrow$ GPIO34 (ADC1);
    \item Botão $\rightarrow$ GPIO15 e GND (usando o \emph{pull-up} interno, como no \cref{cap:gpio}).
\end{itemize}

\section{Projetando a estrutura do estado}

Antes de escrever o laço principal, vale modelar o \textbf{estado} do sistema com os tipos vistos no \cref{cap:structs} — uma prática de organização que se torna cada vez mais valiosa à medida que um firmware cresce:

\begin{codigoc}{Modelando o estado do sistema}
enum ModoOperacao {
    MODO_NORMAL,
    MODO_VERBOSO,
};

struct EstadoEstacao {
    ModoOperacao modo;
    int ultima_leitura;
    float media_movel;
    bool alerta_ativo;
};
\end{codigoc}

Agrupar essas variáveis em uma única \code{struct}, em vez de espalhá-las como variáveis globais soltas, deixa explícito o que compõe ``o estado do sistema'' — e facilita, por exemplo, passar esse estado inteiro para uma função de log, caso o projeto cresça no futuro.

\section{Firmware completo}

\begin{codigoc}{src/main.cpp — estação de monitoramento IoT}
#define PINO_LED    2
#define PINO_LDR    34
#define PINO_BOTAO  15

#define TAMANHO_HISTORICO   10
#define LIMIAR_ESCURO       25.0f   // abaixo disso, ambiente "escuro" (%)
#define INTERVALO_LEITURA   500UL   // ms entre leituras do sensor

enum ModoOperacao {
    MODO_NORMAL,
    MODO_VERBOSO,
};

struct EstadoEstacao {
    ModoOperacao modo;
    int ultima_leitura;
    float media_movel;
    bool alerta_ativo;
};

EstadoEstacao estado = {MODO_NORMAL, 0, 0.0f, false};

int historico[TAMANHO_HISTORICO];
int indice_historico = 0;

volatile bool evento_botao = false;
unsigned long ultima_leitura_ms = 0;

void IRAM_ATTR botao_isr() {
    evento_botao = true;
}

float calcular_media(void) {
    long soma = 0;
    for (int i = 0; i < TAMANHO_HISTORICO; i++) {
        soma += historico[i];
    }
    return (float) soma / TAMANHO_HISTORICO;
}

void setup() {
    Serial.begin(115200);

    pinMode(PINO_LED, OUTPUT);
    pinMode(PINO_BOTAO, INPUT_PULLUP);
    attachInterrupt(digitalPinToInterrupt(PINO_BOTAO), botao_isr, FALLING);

    for (int i = 0; i < TAMANHO_HISTORICO; i++) {
        historico[i] = 0;
    }

    Serial.println("Estacao de monitoramento iniciada em modo NORMAL.");
}

void loop() {
    // 1) verifica se o botao foi pressionado (sem bloquear o laco)
    if (evento_botao) {
        evento_botao = false;
        estado.modo = (estado.modo == MODO_NORMAL) ? MODO_VERBOSO : MODO_NORMAL;

        Serial.print("Modo alterado para: ");
        Serial.println(estado.modo == MODO_NORMAL ? "NORMAL" : "VERBOSO");
    }

    // 2) le o sensor periodicamente, sem usar delay()
    unsigned long agora = millis();
    if (agora - ultima_leitura_ms >= INTERVALO_LEITURA) {
        ultima_leitura_ms = agora;

        estado.ultima_leitura = analogRead(PINO_LDR);
        historico[indice_historico] = estado.ultima_leitura;
        indice_historico = (indice_historico + 1) % TAMANHO_HISTORICO;

        estado.media_movel = calcular_media();
        float luminosidade_pct = (estado.media_movel / 4095.0f) * 100.0f;
        estado.alerta_ativo = (luminosidade_pct < LIMIAR_ESCURO);

        digitalWrite(PINO_LED, estado.alerta_ativo);

        if (estado.modo == MODO_VERBOSO) {
            Serial.print("[LEITURA] bruta=");
            Serial.print(estado.ultima_leitura);
            Serial.print(" media=");
            Serial.print(estado.media_movel);
            Serial.print(" luminosidade=");
            Serial.print(luminosidade_pct);
            Serial.println("%");
        }

        if (estado.alerta_ativo) {
            Serial.println("[ALERTA] Ambiente escuro detectado!");
        }
    }
}
\end{codigoc}

\section{Como o projeto se encaixa}

\begin{itemize}
    \item A \code{struct EstadoEstacao} e a \code{enum ModoOperacao} (\cref{cap:structs}) organizam o estado do sistema de forma clara, em vez de espalhá-lo em variáveis globais soltas;
    \item O vetor \code{historico} (\cref{cap:arranjos}) implementa um \textbf{buffer circular} para a média móvel, suavizando o ruído natural do sensor, exatamente como praticado no \cref{cap:adc};
    \item A leitura do LDR usa a \textbf{ADC} (\cref{cap:adc}) para medir luminosidade de forma contínua, convertida de valor bruto para porcentagem com um \emph{cast} (\cref{cap:tipos});
    \item O botão é lido por \textbf{interrupção} (\cref{cap:interrupcoes}), sinalizando o laço principal por meio de uma variável \code{volatile}, sem nunca bloquear a leitura contínua do sensor;
    \item A leitura periódica do sensor usa a técnica de \textbf{temporização por \code{millis()}} (\cref{cap:interrupcoes}), mantendo \code{loop()} sempre livre para reagir ao botão a qualquer momento;
    \item Todo o comportamento é registrado pela \textbf{porta serial} (\cref{cap:serial}) — mensagens de leitura apenas em modo verboso, e alertas sempre, independentemente do modo.
\end{itemize}

\section{Para continuar praticando}

Esta apostila termina aqui, mas o aprendizado de sistemas embarcados é, acima de tudo, prático. Como próximos passos a partir deste projeto, considere:

\begin{itemize}
    \item Adicionar um segundo sensor (por exemplo, um sensor de temperatura) e expandir a \code{struct EstadoEstacao} para incluí-lo;
    \item Salvar o \code{modo} escolhido na memória flash (usando a biblioteca \code{Preferences} do Arduino para ESP32), para que a estação lembre a última configuração mesmo após ser desligada;
    \item Explorar a conectividade Wi-Fi do \ESP{} para reportar as leituras a um servidor ou aplicativo remoto — o próximo passo natural de qualquer projeto de IoT;
    \item Reescrever o projeto dividindo-o em múltiplos arquivos \code{.h}/\code{.cpp} (\cref{cap:modularizacao}) — por exemplo, separando a lógica do sensor, do histórico e do botão em módulos independentes.
\end{itemize}

Cada um desses passos reutiliza e aprofunda os mesmos fundamentos de C apresentados nesta apostila — a linguagem muda pouco; o que cresce é a sua familiaridade com o hardware e o ecossistema ao redor dela.
